\subsection{Kohomologie von Sphären}

Wir wenden uns nun einer anderen Invariante zu, die für Sphären und ähnliche Typen einfacher zu berechnen ist als die Homotopiegruppen im vorangegangenen Kapitel.
Diese Invariante wird eine Sequenz abelscher Gruppen sein, die Kohomologie genannt wird.
Als Vorabeit werden wir zunächst Eilenberg-MacLane Räume für \(\Z\) betrachten.

\begin{definition}
  Für \(n\in\Z\) sei
  \begin{align*}
    &K(\Z,n)\colonequiv \|S^n\|_n &\text{falls \(n>0\)}\\
    &K(\Z,0)\colonequiv \Z \\
    &K(\Z,n)\colonequiv \eins &\text{falls \(n<0\)}
  \end{align*}
  \(K(\Z,n)\) heißt \(n\)-ter \begriff{Eilenberg--MacLane--Raum}\index{\(K(\Z,n)\)} oder
  \(n\)-te \begriff{Entschleifung} von \(\Z\).
\end{definition}

\begin{lemma}
  \begin{enumerate}
  \item Für alle \(n\in\Z\) gilt:
    \[
      \Omega K(\Z,n+1)=K(\Z,n)
    \]
  \item \(K(\Z,n)\) ist punktiert und für \(n\geq 0\) \(n\)-abgeschnitten und \((n-1)\)-zusammenhängend.
  \item \(\pi_k(K(\Z,n))\) ist \(\Z\) falls \(k=n\) gilt und sonst \(\eins\).
  \end{enumerate}
\end{lemma}

\begin{beweis}
  \begin{enumerate}
  \item Für \(n\geq 2\) können wir \Cref{freudenthal} verwenden um eine Äquivalenz
    \[
      \|S^n\|_n\simeq \Omega\|S^{n+1}\|_{n+1}
    \]
    zu bekommen. Für \(n<1\) haben wir bereits alle Fälle berechnet, es bleibt also der Fall \(n=1\).
    Das heißt wir wollen \(S^1=\Omega\|S^2\|_2(=\|\Omega S^2\|_1) \) zeigen.
    Dazu verwenden wir die Abbildung \(f\) in der Fasersequenz der Hopf-Faserung:
    \begin{center}
      \begin{tikzcd}
            & \Omega^2 S^3\ar[r] & \Omega^2 S^2 \ar[dll,"\Omega f\circ \_^{-1}"] \\
        \Omega S^1\ar[r] & \Omega S^3\ar[r] & \Omega S^2\ar[dll,"f"] \\
        S^1\ar[r] & S^3\ar[r] & S^2 
      \end{tikzcd}
    \end{center}
    Weil \(\Omega S^3\) und \(\Omega^2 S^3\) jeweils \(0\)-zusammenhängend sind, ist
    \(\|\Omega f\circ \_^{-1}\|_0 \) und damit auch \(\|\Omega f\|_0\) eine Äquivalenz.
    Daraus folgt bereits, dass auch \(\Omega \|f\|_1:\Omega\|S^1\|_1\to \Omega\|\Omega S^2\|_1\) eine Äquivalenz ist,
    mit der Kommutativität des folgenden Diagramms für beliebige \(g:A\to_\ast B\):
    \begin{center}
      \begin{tikzcd}
        \|\Omega A\|_0\ar[d,"\simeq"]\ar[r,"\|\Omega f\|_0"] & \|\Omega B\|_0\ar[d,"\simeq"] \\
        \Omega\|A\|_1\ar[r,"\Omega \|f\|_1"] & \Omega\|B\|_1
      \end{tikzcd}
    \end{center}
    (Die wir hier nicht nachrechnen). Da \(\Omega \|f\|_1\) eine Äquivalenz ist und \(\|S^1\|_1=S^1\) und \(\|\Omega S^2\|_1\) beide \(0\)-zusammenhängend sind, folgt dass auch \(\|f\|_1\) eine Äquivalenz ist.
  \item \(K(\Z,n)\) ist punktiert durch \(|\ast|_n:\|S^n\|_n\) bzw.\ \(0:\Z\) und \(\ast:\eins\).
    \(K(\Z,n)\) ist \((n-1)\)-zusammenhängend, da \(\|S^n\|_n\) \((n-1)\)-zusammenhängend ist.
  \end{enumerate}
\end{beweis}

\begin{definition}
  Sei \(X:\mU\) und \(n\in\Z\). Dann ist
  \[
    H^n(X)\colonequiv H^n(X,\Z)\colonequiv \|X\to K(\Z,n)\|_0
  \]
  die \(n\)-te \begriff{Kohomologiegruppe}\index{\(H^n(X)\)}\index{\(H^n(X,\Z)\)} von \(X\) (mit Koeffizienten in \(\Z\)).
\end{definition}

\begin{beispiel}
  \begin{enumerate}[(a)]
  \item \(H^n(\emptyset)\equiv \|\emptyset\to K(\Z,n)\|_0=\|\eins\|_0=\eins\)
  \item \(H^n(\eins)=\|\eins\to K(\Z,n)\|_0=\|K(\Z,n)\|_0\) das ist \(\Z\) für \(n=0\) und sonst \(\eins\).
  \item \(H^n(\zwei)=\|\zwei\to K(\Z,n)\|_0=\|K(\Z,n)\times K(\Z,n)\|_0=\|K(\Z,n)\|_0\times \|K(\Z,n)\|_0 \). Das ist \(\Z\times\Z\) für \(n=0\) und sonst \(\eins\)
  \item \(H^1(S^1)=\Z\). (Übungsblatt 13)
  \end{enumerate}
\end{beispiel}

\begin{bemerkung}
  Sei \(k\geq 0\), dann ist
  \[
    H^n(X)\equiv \|X\to K(\Z,n)\|_0=\|X\to \Omega^k K(\Z,n+k)\|_0=\|\Omega^k(X\to  K(\Z,n+k))\|_0\equiv\pi_k(X\to  K(\Z,n+k))
  \]
  Kohomologiegruppen sind also stets Homotopiegruppen von beliebig hohem Grad.
\end{bemerkung}

\begin{lemma}
  \begin{enumerate}
  \item \(H^n(X)\) ist stets eine abelsche Gruppe.
  \item Für \(f:X\to Y\) ist durch \(f^\ast\colonequiv \|\_\circ f\|_0:H^n(Y)\to H^n(X)\) ein Gruppenhomorphismus gegeben.
  \item Es gelten \(\id^\ast=\id\) und \((g\circ f)^\ast=f^\ast\circ g^\ast\).
  \end{enumerate}
\end{lemma}

\begin{satz}
  Sei \(
  \begin{tikzcd}
    A& C\ar[l,"f",swap]\ar[r,"g"] & B
  \end{tikzcd}
  \). Dann gibt es eine lange exakte Sequenz von Kohomologiegruppen:
  \begin{center}
    \begin{tikzcd}
      \dots\ar[r] & H^{n-1}(A)\oplus H^n(B)\ar[r,"g^\ast-f^\ast"] & H^{n-1}(C)\ar[dll] \\
      H^n(A\sqcup_C B)\ar[r] & H^n(A)\oplus H^n(B)\ar[r,"g^\ast-f^\ast"] & H^n(C)\ar[dll] \\
      H^{n+1}(A\sqcup_C B)\ar[r] & H^{n+1}(A)\oplus H^{n+1}(B)\ar[r,"g^\ast-f^\ast"] & \dots,
    \end{tikzcd}
  \end{center}
  die \begriff{Mayer-Vietoris-Sequenz}. Wobei wir ``\(\oplus\)'' für das binäre Produkt von Gruppen schreiben und ``\(g^\ast-f^\ast\)'' statt ``\((g^\ast)^{-1}\kon f^\ast\)'', weil alle beteiligten Gruppen abelsch sind. 
\end{satz}

\subsection{Kohomologie endlicher CW-Komplexe}

Wir folgen nun \url{https://arxiv.org/pdf/1802.02191}.

\begin{definition}
  \begin{enumerate}
  \item Ein Typ \(X\) heißt \begriff{endlich}, wenn \(\|X=\{0,\dots,n-1\}\|\), wobei \(\{0,\dots,n-1\}\) die \begriff{standard-endliche Menge} mit \(n\) Elementen bezeichnet, die gegeben ist durch
    \[\{0,\dots,n-1\}\colonequiv \sum_{k:\N}k< n\rlap{.}\]
  \item Eine Aussage \(P\) heißt \begriff{entscheidbar}, wenn \(P\vee \neg P\) gilt.
  \item Eine \begriff{entscheidbare Menge} ist eine Menge \(M\), sodass \(x=_My\) stets eine entscheidbare Aussage ist.
  \end{enumerate}
\end{definition}

\begin{beispiel}
  \begin{enumerate}
  \item \(\N\) ist eine entscheidbare Menge.
  \item Jeder endliche Typ ist eine entscheidbare Menge.
  \end{enumerate}
\end{beispiel}

\begin{bemerkung}
  Für einen endlichen Typ \(X\) ist es im Allgemeinen nicht der Fall, dass für eine Menge \(V\) mit Surjektion \(X\to V\), \(V\) wieder endlicher Typ ist.
  Es ist im Allgemeinen auch nicht der Fall, dass ein \(U\) mit einer Einbettung \(U\to X\) wieder endlich ist.
\end{bemerkung}

\begin{bemerkung}
  Sei \(M\) eine entscheidbare Menge. Dann kann \(f:M\to X\) durch Fallunterscheidung definiert werden:
  Für gegebenes \(x:M\) reicht es jeweils \(f\) für \(y=x\) und für \(y\neq x\) anzugeben.
\end{bemerkung}

\begin{definition}
  Sei \(I\) ein Typ und \(X:I\to \mU\) sodass \(X_i\) für alle \(i:I\) punktiert ist.
  Dann ist der \begriff{Wedge}\index{\(\bigvee_{i:I}X_i\)} \(\bigvee_{i:I}X_i\) gegeben durch den folgenden Pushout:
  \begin{center}
    \begin{tikzcd}
      I\ar[r]\ar[d] & \sum_{i:I}X_i\ar[d] \\
      \eins\ar[r] & \bigvee_{i:I}X_i
    \end{tikzcd}
  \end{center}
  weiter gibt es für \(j:I\) eine Inklusionsabbildung:
  \[
    \iota_j:X_j\to \sum_{i:I}X_i\to \bigvee_{i:I}X_i
  \]
  Falls \(I\) entscheidbare Menge ist, können wir für \(j:I\) eine Abbildung
  \(t_j:\sum_{i:I}X_i \to X_j\) durch Fallunterscheidung definieren und damit eine Projektion:
  \[
    \pi_j:\bigvee_{i:I}X_i \to X_j
  \]
  Im Fall \(I=\zwei\) schreiben wir auch \(X_0\vee X_1\)\index{\(\vee\)} statt \(\bigvee_{i:\zwei}X_i}\).
\end{definition}

\begin{definition}
  \begin{enumerate}
  \item Ein Paar \(A,\alpha\) heißt \begriff{endliches CW-Komplex-Datum}, wenn es \(N:\N\) gibt und \(A_i\) für \(0\leq i \leq N\) endlich ist und weiter \(\alpha_{n+1}:A_{n+1}\times S^n\to X_n\), wobei \(X_n\) gegeben ist durch den Pushout:
    \begin{center}
      \begin{tikzcd}
        A_{n+1}\times S^n\ar[r]\ar[d] & X_n\ar[d]  \\
        A_{n+1}\ar[r] & X_{n+1}
      \end{tikzcd}
    \end{center}
    und \(X_0\colonequiv A_0\).
  \item Ein Typ \(X\) ist eine \begriff{endlicher CW-Komplex}, wenn es ein endliches CW-Komplex-Datum \((A,\alpha)\) gibt, sodass \(X=X_N\) gilt mit der Notation aus (a).
  \end{enumerate}
\end{definition}

\begin{beispiel}
  \begin{enumerate}
  \item Endliche Typen sind endliche CW-Komplexe.
  \item \(S^n\) ist endlicher CW-Komplex, gegeben durch
    \[
      A_0\colonequiv \eins,A_1\colonequiv \emptyset,\dots,A_{n-1}\colonequiv \emptyset,A_n\colonequiv \eins
    \]
    zusammen mit den eindeutig bestimmten Abbildungen \(\alpha\).
  \item Der Torus, \(T^2\) ist gegeben durch
    \[
      A_0\colonequiv\eins, A_1\colonequiv \zwei, A_2\colonequiv \eins
    \]
    und der Abbildung \(\alpha_{2}:A_2\times S^1\to X_2=S^1\vee S^1\), die \(l:\ast=_{S^1}\ast\) auf
    \(aba^{-1}b^{-1}\) abbildet, wobei \(a\colonequiv \iota_1(l)\) und \(b\colonequiv \iota_2(l)\).
  \end{enumerate}
\end{beispiel}
